
\section{Background on the domain of inquiry and the data}\label{sec:Background}

\subsection{Adverbials of immediate posteriority in French and German}\label{sec:IPs}

In earlier work we studied ways of expressing immediate posteriority
in several languages, including adverbials such as
Engl. \textit{immediately}, Germ. \textit{sofort} and \mbox{\textit{gleich}}, and
Fr. \textit{immédiatement} and \textit{tout de suite}
(\citealt{Atayanetal2019,GastAtayan2019,Gastetal2019,Atayan2021b,Atayan2021a}).
We will use the abbreviation \qual{\ipadv} in the following.  \ipadv s
locate an event $e$ relative to some reference point $R$ in such a way
that $e$ follows $R$ without a relevant temporal gap.  For example, in
({\ref{ex1}}) Jane's arriving at home ($R$) and her going to bed ($e$)
are temporally adjacent, which is here expressed by
\textit{immediately}.

\begin{exe}
  \ex\label{ex1} When [$_R$ Jane arrived at home] [$_e$ she
  \textit{immediately} went to bed].
\end{exe}
%
Adverbials like \textit{soon} (cf. (\ref{soonex})), \textit{bientôt} (French) and \textit{bald} (German)
are related to \ipadv s insofar as their use implies a relatively
short interval between events or points in time, but the interval
is longer than in the case of \ipadv s.
%

\begin{exe}
  \ex\label{soonex} \textit{Soon} after [$_R$ she arrived], [$_e$ Jane went to bed].
\end{exe}
%
We will refer to the class of adverbials covering both those of
immediate posteriority and \textit{soon}/\textit{bald}/\textit{bientôt}
as `adverbials of proximal posteriority', or \pradv s for
short. \ipadv s are regarded as a special case of \pradv s.\footnote{Note that
arguably, there are also differences in the lengths of the intervals
between, for instance, Germ. \textit{sofort} (no or very short interval) and
\textit{gleich} (somewhat longer interval). We disregard these
differences in the following, as they are not of primary interest to
the questions addressed in our projects.}


Expressions of proximal posteriority tend to be context-sensitive,
insofar as their use is (partially) determined by specific
properties of the context.  In a (simplifying) model of language
production we can assume that there is a given function $f$ that a
speaker wants to express -- say, immediate posteriority -- in some
context $C$ with context variables $\{C_1 \ldots C_m\}$, each of which
comprises a certain number of levels (e.g. `directive', `commissive'
and 'representative' for the variable \textsc{illocutionary force}).
Moreover, there is a set of expression candidates
$W=\{w_1 \ldots w_n\}$ which are compatible with function $f$. One
aspect of use that we are interested in is the probablities
with which the candidates $W=\{w_1 \ldots w_n\}$ are used
in context $C$ to express function $f$.
These probabilities can be regarded as conditional probabilities of
the form $p(w|f,C)$: the probability of using expression $w$, given that we
want to express the function $f$ in context $C$.

To illustrate some relevant context parameters determining the
distribution of \ipadv s, let us consider some examples. \ipadv s may
be used in utterances with representative illocutionary force, as in
(\ref{ex1}) above, or with directive illocutionary force, as in
(\ref{ex2}) below; they can be used anaphorically, with another event
$e$ as a reference point (cf. (\ref{ex1})), or deictically, with the
moment of speech as a reference point (cf. (\ref{ex2})); and they can
be used with intentional (controlled) actions, as in (\ref{ex1}) and
(\ref{ex2}), or with unintentional/spontaneous events, as in
(\ref{ex3}).

\begin{exe}
  \ex\label{ex2} [$_e$ Get out] \textit{immediately}! (directive, deictic, intentional; implicit $R$)    
  \ex\label{ex3} [$_e$ The house of cards  collapsed] \textit{immediately} when $[_R$ I touched it].\\(representative, anaphoric, spontaneous)
\end{exe}
%
\ipadv s often exhibit particularities in their distribution amounting
to interactions between their distributional properties. For example,
Engl. \textit{immediately}, while being compatible with both deictic
reference points (as in \bv{ex2}) and unintentional actions (as in
\bv{ex3}), is not normally used when these two context parameters
co-occur, as in \bv{ex4}, in a reading that is both spontaneous and deictic. This
sentence is not normally understood as a warning about an imminent,
spontaneous collapse, in which case we might utter something like \bv{ex4b}; rather,
\bv{ex4} suggests that the collapse will follow some other event
triggering it, as in (\ref{ex3}) above and \bv{ex5} below.

\begin{exe}
  \ex
  \begin{xlist}
  \ex\label{ex4} Be careful, it will collapse \textit{immediately}!
  \ex\label{ex4b} Be careful, it will collapse \textit{any moment}! (deictic)
  \ex\label{ex5} Be careful, it will collapse \textit{immediately} if you touch
  it! (anaphoric)
\end{xlist}
\end{exe}
%
In translation and interpreting, there are, obviously, additional
factors determining the probability of a specific form being used to
express some function $f$ in a given context $C$. First, it is
well-known that translated (or, more generally speaking, mediated)
language production differs systematically from monolingual language
production (cf. for instance the well-known translation universals
investigated in the tradition of
\citealt{Baker1993,Baker1995}). Second, there are interference
effects, i.e., the choice of a target language form is partially
determined by the form found in the input (often leading to ``shining
through'' effects, cf. \citealt{Teich2003}). There is, thus, an
additional variable to be considered. The probability of using a given
expressions is conditioned by both the context (type) $C$ and the
source language expression $s$. We are thus dealing with a
conditional probability of the form $p(w|f,C,s)$: the probability of
using expression $w$, given that we want to express function $f$ in
context $C$, with expression $s$ in the source text. For example, if
we translate \bv{ex5} into German, we can determine probabilities for
using \textit{gleich} or \textit{sofort} to express immediate
posteriority in the specific context, given that the English original
contains \textit{immediately}. This probability may differ if the
source sentence contained a different adverbial with a similar
meaning, e.g. \textit{right away}.

\subsection{The data}\label{sec:data}

For illustration we use two data sets, both of them from the Europarl corpus
(\citealt{Koehn2005,CartoniMeyer2012,Cartonietal2013}). The first data
set was created in the context of a joint project on \ipadv s carried
out at the Universities of Heidelberg and Jena. The main results of
this project have been published in \citet{Atayanetal2019},
\citet{Gastetal2019}, \citet{GastAtayan2019}, \citet{Atayan2021b} and
\citet{Atayan2021a}, to which the reader is refered for more
information (e.g. concerning the annotation procedure). The data used
for this chapter is based on two originally distinct data sets: (i)
508 annotated observations from German and French in which the \ipadv
s \textit{gleich}, \textit{sofort}, \textit{immédiatement} and
\textit{tout de suite} occur either in originals or in translations,
annotated by a team of annotators as mentioned in the aforementioned
publications;
and (ii) 100 observations each of the two adverbials of proximal
posteriority \textit{bientôt} (French) and \textit{bald} (German),
annotated by V. Atayan.  The data as well as all R- and Python-scripts
used for the analyses and visualizations are available in the
OSF-repository https://osf.io/nufpv/.

%
The data were annotated for the following variables (the specific levels of the variables will be briefly explained when examples are discussed):

\begin{itemize}[noitemsep]
\item Adverbial in original and translation %(as observed in the data)
\item French and German tense and mood %(as observed in the data)
\item Illocutionary force/modality%\footnote{The original annotation categories were: `nicht-modal/assertiv', `Aufforderung/Wunsch', `Versprechen/Drohung', `Vermutung', `keine Modalitaetsmarkierung', `Moeglichkeit/Faehigkeit', `sachliche Notwendigkeit', `Warnung'. For the analysis, some of these levels were merged.}
  % \begin{compactitem}
  % \item comm[issive] (promise)
  % \item dir[ective].boul[etic] (instruction or wish)
  % \item repr[esentative] (description of the world)
  % \item poss[ibility] (modality)
  % \item nec[essity] (modality)
  % \item n[o].illoc[utionary force] (e.g. non-finite predicates)
  % \end{compactitem}
\item Presence of a syntactic modifier% (modified, unmodified)
  % \begin{compactitem}
  % \item modified
  % \item unmodified
  % \end{compactitem}
\item Intentionality of the predicate% (intentional, non-intentional)
  % \begin{compactitem}
  % \item intentional
  % \item non-intentional
  % \end{compactitem}
\item Agent of the main predicate%\footnote{The original annotation categories were: `1sg', `3sg', `1pl', `3pl', `2pl', `2sg', `3-belebt-unpezifisch', `3-unbelebt', `3-belebt-unpersoenlich'. For the analysis, some of these levels were neutralized.\label{AgentNote}}
  % \begin{compactitem}
  % \item 1sg (1. person sg.)
  % \item 1pl (1. person pl.)
  % \item 3ps (2. person)
  % \item 3ps (3. person)
  % \end{compactitem}
\item Reference point%\footnote{The original levels were deic[tic]-dir[ect] (moment of speech is reference point); deic[tic]-phant[asma] (deictic transposition); chron[ological]-ev[ent] (event mentioned in context is reference  point); chronological]-t[i]m[e] (point in time mentioned in  context is reference point).} %with levels
  % \begin{compactitem}%[noitemsep]
  % \item chron[ological]
  % \item deic[tic]
  % \end{compactitem}
\end{itemize}
%


%%% Local Variables:
%%% mode: latex
%%% TeX-master: "main"
%%% End:
